% =====================================================================
%  VORLAGE LABORBERICHT (MCI, Klasse mcidoc, state=Report)
%  Vorgaben: siehe CLAUDE.md (MCI-Formalkriterien, M-Codes).
%  Kompilieren: latexmk -pdf main.tex (MiKTeX, biber)
%  Kopiervorlagen für Tabellen, Abbildungen, Gleichungen, Zitate: Vorlagen.tex
% =====================================================================
\documentclass[state=Report]{mcidoc}
\usepackage{tabularx}
\usepackage{booktabs}                             % Tabellen ohne vertikale Linien (M057)
\usepackage{pgfplots}                             % Diagramme direkt in LaTeX
\pgfplotsset{compat=1.18}

% ---------- Ergänzungen für die MCI-Formalkriterien ----------
\sisetup{locale = UK, uncertainty-mode = separate} % Dezimalpunkt (M030), Unsicherheit mit ± (M031)
\KOMAoptions{listof=totoc}                        % Verzeichnisse erscheinen im Inhaltsverzeichnis
\KOMAoptions{numbers=noendperiod}                 % sonst "Figure 2.1.:" sobald ein Anhang existiert (M018)
\AtBeginEnvironment{table}{\setstretch{1.1}}      % Tabellen kompakter als Fließtext (1.5), einheitlich (M018)
\newenvironment{legendtable}[1]                  % Legende zu Gleichungen (M042), gleicher Zeilenabstand wie Tabellen
  {\par\setstretch{1.1}\noindent\begin{tabular}{#1}}
  {\end{tabular}\par}
\usepackage[autostyle=true]{csquotes}            % englische Anführungszeichen
\usepackage[backend=biber, style=ieee]{biblatex}  % IEEE-Zitierstil (M080)
\addbibresource{references.bib}
\newcounter{romanpage}
% Seitenzahl + Name auch auf Kapitel-Startseiten (Stil "plain"), sonst fehlen sie dort (M010)
\makeatletter
\ofoot*{\thepage}
\ifoot*{\Global@StudentName}
\makeatother
\KOMAoptions{plainfootsepline=true}

% ---------- Metadaten: ANPASSEN ----------
\SetMCI{Language = EN,                 % EN oder DE
    Degree = BSc,
    Title = {Lab Report},
    StudentID = {00000000, 00000000},  % Matrikelnummern, gleiche Reihenfolge wie Namen
    Student/Name = {First Author, Second Author},
    Report/Supervisor/Name = {Supervisor Name},
    Report/Cohort = {BA-XXXX-00-0X},
    Report/Lecture = {Course Name},
    Report/Date = {January 1, 2026},   % Labortermin bzw. Abgabedatum laut Anleitung
    Report/Mode = Lab,                 % Lab -> "Lab Supervisor", Assignment -> "Supervisor"
    Report/ToC/Primary = unset,        % Verzeichnisse werden unten selbst gesetzt (M010)
    Report/ToC/Figures = unset,
    Report/ToC/Tables  = unset,
}

\begin{document}

% ---------- Vorspann: römische Seitenzahlen (M010) ----------
\pagenumbering{Roman}
\tableofcontents
\clearpage
\setcounter{romanpage}{\value{page}}
\pagenumbering{arabic}

% ---------- Hauptteil: arabische Seitenzahlen ----------
\chapter{Introduction}\label{ch:introduction}
% Passiv + Präsens, keine Ich-/Wir-Form (M019). Quellen direkt am Satz zitieren (M070).

\section{General Description}\label{sec:general}
% Was wird untersucht? Aufbau, Messgröße, verwendete Hard- und Software in 3--5 Sätzen.

\section{Motivation}\label{sec:motivation}
% Warum ist das Thema relevant? Anwendungsbeispiele, Bedarf an Modell/Messung.

\section{Learning Goal}\label{sec:goal}
% Lernziele laut Laboranleitung in eigenen Worten (keine Aufzählung 1:1 übernehmen, M104).

\chapter{Procedure and Method}\label{ch:procedure}

\section{Experimental Setup}\label{sec:setup}
% Materialliste + beschriftete Versuchsskizze (M086), Hersteller/Typ/Präzision (M088).
The equipment is listed in Table~\ref{tab:equipment}.

\begin{table}[H]
  \centering
  \caption[List of equipment]{List of equipment used for the experiment}
  \label{tab:equipment}
  \begin{tabularx}{\textwidth}{l >{\raggedright\arraybackslash}X l}
    \toprule
    Component & Description & Type \\
    \midrule
    Device      & description of the device & Manufacturer Type \\
    Software    & purpose of the software   & Name Version \\
    \bottomrule
  \end{tabularx}
\end{table}

\section{Method}\label{sec:method}
% Vorgehen, Modelle und Grundgleichungen. Zentrale Gleichungen mit Legende (M042),
% nur referenzierte Gleichungen nummerieren (equation), sonst equation* (M040).

\chapter{Results}\label{ch:results}
% Nur eigene Messungen und Ergebnisse. Beispielrechnung je Rechenschritt mit allen
% Einheiten und Bezug auf die Messdaten (M099--M101). Diagramme: Vorlagen.tex, Kapitel 5.
% Messdaten für pgfplots als CSV in data/ ablegen.

\section{Measurement Results}\label{sec:res-measurement}

\chapter{Discussion and Interpretation}\label{ch:discussion}
% Meist der wichtigste Teil: hier ausreichend Platz einplanen.

\section{Comparison of Measurement, Theory and Simulation}\label{sec:comparison}
% Abweichungen benennen und begründen.

\section{Limitations}\label{sec:limitations}
% Messunsicherheit, Reproduzierbarkeit, Modellannahmen.

\section{Conclusion}\label{sec:conclusion}
% Zusammenfassung der wichtigsten Erkenntnisse und Ausblick.


% ---------- Nachspann: römisch fortlaufend (M010) ----------
\clearpage
\pagenumbering{Roman}
\setcounter{page}{\value{romanpage}}

\printbibliography[heading=bibintoc]
\listoffigures
\listoftables

\addchap{List of Symbols}              % Formelzeichenverzeichnis (M049): jedes Formelzeichen eintragen
\begingroup
\setlength{\parfillskip}{0pt plus 1fil}  % sonst 10 pt breiter als der Satzspiegel (M057)
\noindent
\begin{tabularx}{\textwidth}{l X >{\raggedright\arraybackslash}p{0.2\textwidth}}
  \toprule
  Symbol & Description & Unit \\
  \midrule
  $t$            & time & \unit{\second} \\
  $u$            & voltage & \unit{\volt} \\
  \bottomrule
\end{tabularx}
\par
\endgroup

% ---------- Anhang: römisch fortlaufend, zählt meist nicht zur Seitenvorgabe ----------
% Auskommentiert, solange er leer ist; zum Aktivieren beide Zeilen einkommentieren.
% \appendix
% % =====================================================================
% Anhang: Seitenzahlen laufen römisch weiter (M010), Kapitel werden A, B, C ... nummeriert.
% Jeder Inhalt hier muss im Haupttext referenziert werden (M053, M060), z. B.:
%   ... see Figure~\ref{fig:app-example} in Appendix~\ref{app:data}.
% Zwei Abbildungen nebeneinander: floatrow + \ffigbox (siehe Vorlagen.tex, Kapitel 5).
% =====================================================================
\chapter{Additional Data}\label{app:data}


\end{document}
